\begin{objectivebox}[이 장에서 붙잡을 질문]
같은 입력에서 같은 출력을 내는 회로가 여러 개라면, 무엇이 같고 어떤 비용이 달라질까?
\end{objectivebox}

\section{둘 가운데 정확히 하나만 고르는 규칙}

두 개의 벽 스위치 $A$와 $B$가 있다고 하자. 두 스위치의 상태가 서로 다를 때만 표시등 $Y$를 켜고 싶다. 둘 다 꺼져 있거나 둘 다 켜져 있으면 표시등은 꺼진다. 이 규칙의 네 경우를 빠짐없이 적으면 다음과 같다.

\begin{center}
\small
\begin{tabular}{ccc}
\toprule
\sffamily\bfseries $A$ & \sffamily\bfseries $B$ & \sffamily\bfseries $Y$: 표시등 \\
\midrule
0 & 0 & 0 \\
0 & 1 & 1 \\
1 & 0 & 1 \\
1 & 1 & 0 \\
\bottomrule
\end{tabular}
\end{center}

두 입력 가운데 \emph{정확히 하나만} $1$일 때 출력이 $1$인 함수를 \keyterm{XOR}(exclusive OR, 배타적 논리합)라고 한다. XOR의 마지막 열은 $0,1,1,0$이다. 1장의 AND는 두 입력이 모두 $1$일 때만 $1$이었지만, XOR는 둘이 서로 다를 때만 $1$이다.

진리표는 무엇을 얻어야 하는지를 이미 완전하게 정했다. 그러나 NAND 게이트를 어떻게 연결해야 하는지는 말하지 않는다. 이제 같은 마지막 열을 만드는 두 경로를 각각 끝까지 계산해 보자.

\section{방법 A: 문장을 작은 기능으로 나눈다}

``정확히 하나만 $1$''이라는 말을 두 조건으로 나눌 수 있다.

\begin{enumerate}
  \item 둘 가운데 하나 이상은 $1$이다.
  \item 둘 다 $1$인 경우는 제외한다.
\end{enumerate}

첫째 조건의 ``하나 이상''을 OR(inclusive OR, 논리합)라고 한다. OR는 $A$와 $B$ 가운데 적어도 하나가 $1$이면 $1$이다. 둘 다 $1$이어도 $1$이라는 점이 XOR와 다르다. 여기서는 OR를 새 핵심어로 외우기보다 일상어 ``또는''을 옮긴 중간 기능으로만 사용한다. 둘째 조건은 AND의 결과를 뒤집은 NAND가 바로 알려 준다. 따라서 두 조건을 함께 만족시키면 XOR가 된다.

이처럼 큰 규칙을 작은 함수와 그 연결 관계로 적은 것을 여기서는 \keyterm{논리식}(logic expression)이라고 부르자. 논리식은 진리표와 다른 방식으로 같은 함수를 설명한다. 진리표는 모든 입력을 행으로 펼치고, 논리식은 중간 기능을 어떤 순서로 조합하는지 보여 준다.

1장에서 배운 NAND만으로 이 문장을 옮기면 다음 여섯 단계를 얻는다.

\begin{align*}
  N_A &= \operatorname{NAND}(A,A) &&= \operatorname{NOT}(A),\\
  N_B &= \operatorname{NAND}(B,B) &&= \operatorname{NOT}(B),\\
  O   &= \operatorname{NAND}(N_A,N_B) &&= A\ \operatorname{OR}\ B,\\
  N   &= \operatorname{NAND}(A,B) &&= \operatorname{NOT}(A\ \operatorname{AND}\ B),\\
  T   &= \operatorname{NAND}(O,N) &&= \operatorname{NOT}(A\ \operatorname{XOR}\ B),\\
  Y   &= \operatorname{NAND}(T,T) &&= A\ \operatorname{XOR}\ B.
\end{align*}

$N_A$와 $N_B$는 각각 입력을 뒤집는다. 두 뒤집힌 값을 NAND에 넣으면 $A$ 또는 $B$ 가운데 하나 이상이 $1$인 $O$를 얻는다. $N$은 두 입력이 모두 $1$인 경우에만 $0$이다. $O$와 $N$이 모두 $1$인 경우가 바로 XOR가 $1$이어야 하는 두 경우다. 다섯째 NAND는 그 결과를 한 번 뒤집은 $T$를 만들고, 마지막 NAND가 다시 뒤집어 $Y$를 되찾는다.

말만으로는 중간값 하나가 틀려도 알아차리기 어렵다. 네 입력에서 여섯 값을 모두 계산하면 다음과 같다.

\begin{center}
\scriptsize
\setlength{\tabcolsep}{4.2pt}
\begin{tabular}{cccccccc}
\toprule
\sffamily\bfseries $A$ & \sffamily\bfseries $B$ & \sffamily\bfseries $N_A$ & \sffamily\bfseries $N_B$ & \sffamily\bfseries $O$ & \sffamily\bfseries $N$ & \sffamily\bfseries $T$ & \sffamily\bfseries $Y$ \\
\midrule
0 & 0 & 1 & 1 & 0 & 1 & 1 & 0 \\
0 & 1 & 1 & 0 & 1 & 1 & 0 & 1 \\
1 & 0 & 0 & 1 & 1 & 1 & 0 & 1 \\
1 & 1 & 0 & 0 & 1 & 0 & 1 & 0 \\
\bottomrule
\end{tabular}
\end{center}

마지막 열은 목표 진리표의 $0,1,1,0$과 같다. 방법 A는 NAND 게이트 여섯 개로 XOR를 구현한다.

\section{방법 B: 같은 중간값을 더 알뜰하게 쓴다}

문장을 충실히 옮긴 회로가 유일한 답은 아니다. 첫 NAND의 결과를 두 갈래에서 다시 쓰면 더 짧은 분해를 만들 수 있다.

\begin{align*}
  N_1 &= \operatorname{NAND}(A,B),\\
  N_2 &= \operatorname{NAND}(A,N_1),\\
  N_3 &= \operatorname{NAND}(B,N_1),\\
  Y   &= \operatorname{NAND}(N_2,N_3).
\end{align*}

여기서는 $N_1$을 $N_2$와 $N_3$이 함께 사용한다. 식의 모양만 보고 XOR라고 믿지 말고 네 입력을 다시 계산하자.

\begin{center}
\small
\begin{tabular}{cccccc}
\toprule
\sffamily\bfseries $A$ & \sffamily\bfseries $B$ & \sffamily\bfseries $N_1$ & \sffamily\bfseries $N_2$ & \sffamily\bfseries $N_3$ & \sffamily\bfseries $Y$ \\
\midrule
0 & 0 & 1 & 1 & 1 & 0 \\
0 & 1 & 1 & 1 & 0 & 1 \\
1 & 0 & 1 & 0 & 1 & 1 \\
1 & 1 & 0 & 1 & 1 & 0 \\
\bottomrule
\end{tabular}
\end{center}

방법 B의 마지막 열도 $0,1,1,0$이다. 방법 A와 중간값은 다르지만 입력마다 최종 출력은 같다.

\section{같은 함수라는 말은 무엇을 보존하는가}

두 논리식이 가능한 모든 입력에서 같은 출력을 내면 두 식은 \keyterm{동치}(equivalent)라고 한다. 여기서는 입력이 두 개뿐이므로 네 행을 전부 대조해 동치를 확인했다. 중간값의 이름과 순서가 달라도 마지막 열이 같으면 두 회로는 같은 XOR 함수를 구현한다.

동치는 구현의 모든 성질이 같다는 뜻이 아니다. 같은 것은 $A,B$를 넣었을 때 나오는 $Y$다. 게이트의 개수, 가장 긴 의존 경로, 중간 신호가 갈라지는 횟수, 배선의 위치는 달라질 수 있다. 수학적 기능을 보존하면서 물리적 비용을 바꿀 여지가 바로 이 차이에서 생긴다.

NAND는 조합만으로 NOT, AND, OR를 만들 수 있고, 이 기능들을 다시 조합하면 유한한 입력의 모든 불 함수를 만들 수 있다. 이런 성질을 \keyterm{보편 게이트}(universal gate)라고 한다\parencite{mit6004truth,nand2tetrisproject1}. NAND가 유일한 보편 게이트라는 뜻은 아니다. NOR도 혼자서 같은 역할을 할 수 있다. 또한 논리적으로 NAND만으로 만들 수 있다는 사실이 실제 칩을 NAND 셀로만 만드는 것이 언제나 가장 좋은 선택이라는 뜻도 아니다.

\begin{bookfigure}{하나의 XOR 진리표에서 두 NAND 분해가 갈라진다}{fig:ch02-one-function-two-decompositions}
\begin{tikzpicture}[diagram]
  \path[use as bounding box] (0,0) rectangle (13.0,8.15);

  \draw[diagram panel,fill=black!5] (5.18,1.02) rectangle (7.82,7.95);
  \node[diagram panel title,anchor=north] at (6.50,7.63) {보존할 XOR 명세};
  \node[concept box,minimum width=7mm] at (5.72,6.65) {$A$};
  \node[concept box,minimum width=7mm] at (6.50,6.65) {$B$};
  \node[concept emphasis,minimum width=7mm] at (7.28,6.65) {$Y$};
  \foreach \y/\a/\b/\out in {
    5.89/0/0/0,
    5.20/0/1/1,
    4.51/1/0/1,
    3.82/1/1/0
  }{
    \node[gate,minimum width=7mm] at (5.72,\y) {\a};
    \node[gate,minimum width=7mm] at (6.50,\y) {\b};
    \node[path gate,minimum width=7mm] at (7.28,\y) {\out};
  }
  \node[diagram note,align=center] at (6.50,2.76) {같은 마지막 열\\$0,1,1,0$};
  \node[diagram note,align=center,text=black!62] at (6.50,1.55) {무엇을\\해야 하는가};

  \draw[concept dashed arrow] (5.00,4.52) -- (4.52,4.52);
  \draw[concept dashed arrow] (8.00,4.52) -- (8.48,4.52);

  \draw[diagram panel] (0.05,1.02) rectangle (4.34,7.95);
  \node[diagram panel title,anchor=north] at (2.20,7.63) {방법 A · 문장을 따라간다};
  \node[diagram note,text=black!62] at (0.69,6.82) {1층};
  \node[diagram note,text=black!62] at (1.69,6.82) {2층};
  \node[diagram note,text=black!62] at (2.69,6.82) {3층};
  \node[diagram note,text=black!62] at (3.69,6.82) {4층};
  \foreach \x in {1.19,2.19,3.19}{\draw[black!18,dashed] (\x,2.00)--(\x,6.45);}

  \node[gate,minimum width=8mm] (na) at (0.69,5.80) {$N_A$};
  \node[gate,minimum width=8mm] (nb) at (0.69,4.64) {$N_B$};
  \node[gate,minimum width=8mm] (n)  at (0.69,3.48) {$N$};
  \node[path gate,minimum width=8mm] (o) at (1.69,5.22) {$O$};
  \node[path gate,minimum width=8mm] (t) at (2.69,4.64) {$T$};
  \node[path gate,minimum width=8mm] (ya) at (3.69,4.64) {$Y$};
  \draw[path wire,shorten <=0.3mm,shorten >=0.8mm] (na.east) -- (o.north west);
  \draw[path wire,shorten <=0.3mm,shorten >=0.8mm] (nb.east) -- (o.south west);
  \draw[path wire,shorten <=0.3mm,shorten >=0.8mm] (o.east) -- (t.north west);
  \draw[path wire,shorten <=0.3mm,shorten >=0.8mm] (n.east) --
    node[pos=0.48,below=0.8mm,fill=white,inner sep=0.5pt,
      font=\DiagramBodyFont,text=black!62] {2층 없음}
    (t.south west);
  \draw[path wire,shorten <=0.3mm,shorten >=0.8mm] (t.east) -- (ya.west);
  \node[diagram note,align=center] at (2.20,2.56) {$N_A,N_B,N\rightarrow O\rightarrow T\rightarrow Y$};
  \node[concept emphasis,minimum width=31mm,minimum height=8mm] at (2.20,1.62) {NAND 6개 · 4층};

  \draw[diagram panel] (8.66,1.02) rectangle (12.95,7.95);
  \node[diagram panel title,anchor=north] at (10.80,7.63) {방법 B · 중간값을 다시 쓴다};
  \node[diagram note,text=black!62] at (9.55,6.82) {1층};
  \node[diagram note,text=black!62] at (10.80,6.82) {2층};
  \node[diagram note,text=black!62] at (12.05,6.82) {3층};
  \draw[black!18,dashed] (10.18,2.00)--(10.18,6.45);
  \draw[black!18,dashed] (11.43,2.00)--(11.43,6.45);

  \node[gate,minimum width=10mm] (n1) at (9.55,4.64) {$N_1$};
  \node[path gate,minimum width=10mm] (n2) at (10.80,5.28) {$N_2$};
  \node[path gate,minimum width=10mm] (n3) at (10.80,4.00) {$N_3$};
  \node[path gate,minimum width=10mm] (yb) at (12.05,4.64) {$Y$};
  \draw[path wire,shorten <=0.3mm,shorten >=0.8mm] (n1.east) -- (n2.west);
  \draw[path wire,shorten <=0.3mm,shorten >=0.8mm] (n1.east) -- (n3.west);
  \draw[path wire,shorten <=0.3mm,shorten >=0.8mm] (n2.east) -- (yb.north west);
  \draw[path wire,shorten <=0.3mm,shorten >=0.8mm] (n3.east) -- (yb.south west);
  \node[diagram note,align=center] at (10.80,2.56) {$N_1\rightarrow N_2,N_3\rightarrow Y$};
  \node[concept emphasis,minimum width=31mm,minimum height=8mm] at (10.80,1.62) {NAND 4개 · 3층};

  \draw[gate group] (0.10,0.18) rectangle (12.92,0.78);
  \node[diagram footer] at (6.50,0.48)
    {같은 출력은 보존하고, 중간값과 연결 방법은 바꿀 수 있다};
\end{tikzpicture}
\end{bookfigure}


\figref{fig:ch02-one-function-two-decompositions}은 하나의 XOR 명세에서 두 구현이 갈라지는 모습을 보여 준다. 방법 A는 문장의 두 조건을 드러내므로 왜 XOR가 되는지 따라가기 쉽다. 방법 B는 첫 중간값을 다시 사용해 게이트 수와 의존 경로를 줄인다. 어느 쪽도 진리표가 강제한 유일한 회로는 아니다.

\section{게이트 수와 회로 깊이는 다른 숫자다}

먼저 필요한 NAND 상자를 세어 보자. 방법 A에는 $N_A,N_B,O,N,T,Y$, 모두 여섯 개가 있다. 방법 B에는 $N_1,N_2,N_3,Y$, 모두 네 개가 있다.

다음에는 입력 변화가 출력에 이르기까지 거쳐야 하는 가장 긴 연결을 센다. 이 책에서는 같은 2입력 NAND 한 번을 한 층으로 세고, 가장 많은 NAND를 차례로 통과하는 경로의 층수를 \keyterm{회로 깊이}(circuit depth)라고 부른다.

방법 A에서 가장 긴 경로 하나는

\[
  A \longrightarrow N_A \longrightarrow O \longrightarrow T \longrightarrow Y
\]

이므로 NAND 네 층이다. $B\rightarrow N_B\rightarrow O\rightarrow T\rightarrow Y$도 네 층이다. 반면 방법 B의 가장 긴 경로는

\[
  A \longrightarrow N_1 \longrightarrow N_2 \longrightarrow Y
\]

또는 $B\rightarrow N_1\rightarrow N_3\rightarrow Y$이므로 세 층이다.

게이트 수는 필요한 기능 상자의 총량을 세고, 회로 깊이는 그중 차례로 기다려야 하는 가장 긴 의존 사슬을 센다. 네 게이트가 모두 동시에 작동하는 것도 아니고, 게이트 여섯 개라는 숫자가 곧 여섯 단위의 실행 시간을 뜻하는 것도 아니다. 총량과 가장 긴 경로는 서로 다른 질문에 답한다.

\figref{fig:ch02-same-output-different-cost}의 비교는 의도적으로 제한되어 있다. 같은 2입력 NAND를 한 단위로 센 모형에서는 방법 B가 게이트 두 개 적고 가장 긴 경로도 한 층 짧다. 이 모형은 배선 길이, 입력 부하와 한 신호가 여러 곳으로 갈라지는 \keyterm{팬아웃}(fan-out)을 세지 않으므로 결론도 그 범위 안에 한정된다.

\begin{bookfigure}{같은 출력과 달라진 비용을 따로 센다}{fig:ch02-same-output-different-cost}
\begin{tikzpicture}[diagram]
  \path[use as bounding box] (0,0) rectangle (13.0,5.80);

  \node[concept box,minimum width=31mm,minimum height=10mm] at (1.82,5.15) {구현};
  \node[concept box,minimum width=38mm,minimum height=10mm] at (5.40,5.15) {보존한 XOR 출력};
  \node[concept box,minimum width=25mm,minimum height=10mm] at (8.77,5.15) {NAND 수};
  \node[concept box,minimum width=28mm,minimum height=10mm] at (11.52,5.15) {가장 긴 경로};

  \draw[diagram panel,fill=black!4] (0.05,3.05) rectangle (12.95,4.55);
  \node[diagram panel title,anchor=west] at (0.40,3.80) {방법 A};
  \node[diagram note,align=center] at (5.40,3.80) {$0,1,1,0$};
  \foreach \x in {8.12,8.38,8.64,8.90,9.16,9.42}{\node[gate,minimum width=4.3mm,minimum height=5mm] at (\x,3.80) {};}
  \node[diagram note,anchor=west] at (9.65,3.80) {6개};
  \foreach \x in {10.30,10.82,11.34,11.86}{\node[path gate,minimum width=5.6mm,minimum height=5mm] at (\x,3.80) {};}
  \node[diagram note,anchor=east] at (12.75,3.80) {4층};

  \draw[diagram panel] (0.05,1.32) rectangle (12.95,2.82);
  \node[diagram panel title,anchor=west] at (0.40,2.07) {방법 B};
  \node[diagram note,align=center] at (5.40,2.07) {$0,1,1,0$};
  \foreach \x in {8.38,8.72,9.06,9.40}{\node[gate,minimum width=5.6mm,minimum height=5mm] at (\x,2.07) {};}
  \node[diagram note,anchor=west] at (9.70,2.07) {4개};
  \foreach \x in {10.58,11.24,11.90}{\node[path gate,minimum width=7.0mm,minimum height=5mm] at (\x,2.07) {};}
  \node[diagram note,anchor=east] at (12.75,2.07) {3층};

  \node[concept emphasis,minimum width=101mm,minimum height=9mm] at (6.50,0.66)
    {교육 모형의 결론: 방법 B는 NAND 2개가 적고 가장 긴 경로가 1층 짧다};
\end{tikzpicture}
\end{bookfigure}


그럼에도 이 작은 비교는 중요한 원리를 남긴다. 원하는 함수가 정해져도 분해 방식은 하나로 고정되지 않는다. 같은 답을 지키면서 중간값을 고르고 다시 쓰는 방법에 따라 필요한 자원과 의존 경로가 달라질 수 있다.

다음 장에서는 논리값 $0$과 $1$에 자릿값을 부여한다. 그러면 같은 진리표가 참과 거짓의 규칙을 넘어 한 자리 덧셈을 수행하기 시작한다. 질문도 ``같은 함수를 어떻게 줄일까''에서 ``비트의 위치를 숫자의 크기로 어떻게 읽을까''로 이동한다.

\section*{이번 장의 선택과 비용}
\addcontentsline{toc}{section}{이번 장의 선택과 비용}

\begin{center}
\small
\begin{tabular}{>{\raggedright\arraybackslash\sffamily\bfseries}p{0.29\textwidth} >{\raggedright\arraybackslash}p{0.58\textwidth}}
\toprule
보존한 기능·정확한 결과 & XOR의 네 입력·출력 관계 $0,1,1,0$ \\
\midrule
계산이 요구하는 양 & 네 입력 경우에서 같은 불 함수와 같은 최종 출력 \\
\midrule
기계가 제공할 자원 & 방법 A는 NAND 6개·가장 긴 경로 4층, 방법 B는 NAND 4개·가장 긴 경로 3층 \\
\midrule
설계 대가 & 중간값을 다시 쓰는 분해는 NAND를 2개 줄이고 가장 긴 NAND 경로를 1층 줄임 \\
\midrule
아직 해결되지 않은 제약 & 논리값 $0$과 $1$을 숫자의 자릿값으로 읽고 같은 회로 계층으로 산술을 만드는 방법 \\
\bottomrule
\end{tabular}
\end{center}
